O artigo examina como preços reais de soja e milho se associam à conversão bruta da vegetação nativa para agricultura em Mato Grosso, considerando capacidade produtiva potencial e perdas anteriores. O painel municipal de 2001--2024 combina MapBiomas, PAM, GAEZ, preços internacionais e clima, com efeitos fixos municipais e de bioma por ano. A associação preço×capacidade diminui quando se exige o histórico necessário às perdas. Para +10\% no preço e dez pontos percentuais adicionais de perda, os contrastes são 0,0040 p.p. na soja e $-0,0001$ p.p. no milho, com intervalos de 95\% abrangendo zero nas três inferências. Comparações temporais, territoriais e com pastagem mostram heterogeneidade sensível ao suporte e à inferência. A exposição geográfica potencial à Moratória e a cobertura florestal inicial em torno da reforma de 2012 não identificam efeitos das políticas. Os resultados delimitam associações condicionais, sem demonstrar mecanismos financeiros ou ausência de resposta aos incentivos.

\medskip
\noindent\textbf{Palavras-chave:} conversão da vegetação; preços agrícolas; capacidade produtiva potencial; quebras de safra; Mato Grosso.
